%% kapa: exam=1 round=35 year=2024 period=1 form=A subject=economics subject_ko=경제학원론 seq=34 number=74
%% source: Q-Net 감정평가사 기출문제 (cst00309) artlSeq=5213229 file=raw/1008/5213229/1교시 문제지 원본.pdf
%% conversion: kordoc 4.24.51 -> tools/gen1.py (auto); review=proofread-sonnet; flags=none
\begin{question}{74}
고정환율제를 채택하고 있는 정부가 시장균형환율보다 높은 수준의 환율을 설정했다고 할 때, 즉 자국통화가치를 균형수준보다 낮게 설정한 경우, 옳은 것을 모두 고른 것은?
\begin{framedtext}
ㄱ. 투기적 공격이 발생하면 국내 통화공급이 감소한다. \par
ㄴ. 투기적 공격이 발생하면 외환보유고가 감소한다. \par
ㄷ. 자본이동이 완전히 자유로운 경우, 중앙은행은 독립적으로 통화공급을 결정할 수 없다. \par
ㄹ. 투자자들이 국내통화의 평가절상을 기대하게 되면, 국내통화로 계산된 외국채권의 기대수익률이 하락한다. \par
\end{framedtext}
\begin{choices}
  \item ㄱ, ㄴ
  \item ㄱ, ㄹ
  \item ㄴ, ㄷ
  \item ㄷ, ㄹ
  \item ㄴ, ㄷ, ㄹ
\end{choices}
\end{question}
\answer{4}
